% =============================================================================
% Appendix: Cut Content
% =============================================================================
% SPDX-FileCopyrightText: 2026 Drewan Rutherford <RutherfordD@cardiff.ac.uk>
% SPDX-License-Identifier: LPPL-1.3c
% =============================================================================

This appendix covers the work completed on CodeHeist that did not make it into the Core Systems nor the report.

This appendix has the following sections:

\begin{itemize}
    \item \textbf{Section~\ref{sec:cut/scenarioDesign}} \nameref{sec:cut/scenarioDesign}, which details the planned scenario.
    \item \textbf{Section~\ref{sec:cut/plannedLevels}} \nameref{sec:cut/plannedLevels}, which details an overview of structure of the game.
    \item \textbf{Section~\ref{sec:cut/puzzleIdeas}} \nameref{sec:cut/puzzleIdeas}, which details puzzle ideas.
\end{itemize}

\section{Scenario Design}\label{sec:cut/scenarioDesign}

There were to be two main characters. First, a blank slate, controlled by the player henceforth referred to as the Player. And second, their sister Eve. The name Eve is a reference to the cyber security where Eve is is often use of an malicious actor listening into messages between the placeholders Alice and Bob.

Eve, is trying to steal an initially unknown object, presumably money, from the bank (hence the name of the game). The Player, their younger sibling, is helping them by controlling an autonomous drone CatBot that was made by Eve. As CatBot proceeds further into the depths, more and more stuff goes wrong with the Auto Controllers, as an explanation of why they need to implement basic mechanics they previously had access to.

The story would take on an absurd tone. Whilst not to the same degree as absurd indie games such as \textquote{Turnip Boy Commits Tax Evasion} by \citet{TurnitTaxEvation} or \textquote{Thank Goodness You’re Here!} by \citet{ThankGoodnessURHere}\footnote{The game would also be more appropriate for a younger demographic. This is more representative of tone than the finished result.}, the game was going make use of absurd humour and poke fun at the fact that it was an educational game in a self aware, but distracting way. Most of this was going to come in the form of interactions between the older Eve and the player, potentially similarly to how the player is interacted with in Rhythm Doctor \cite{RhythmDoctor} where characters talk to the player directly using the framing device of the player's remote connection into the hospital, where the player can only respond using in game interaction.

The plot would, beyond the sillier heist elements, would focus on the relationship between Eve and the Player. When the game takes place, Eve had left University and had changed significantly in an unspecified way, and the game is about them and their sibling reconnecting whilst staging the heist.

As one of the later turning points in the story (the exact point was not determined, though it was expected to be around the end of the second act, if the three act structure was used), Eve's motivations would be revealed. It would turn out that she is stealing the launch codes for the orbital death laser to destroy said codes as revenge against her former super villain boss from an unpaid internship, Bob for being mean to her partner. This was meant to be where the more mundane characters intersects with the absurdity of the Bank the itself.

This scenario was chosen because it mirrors relationships that many people, especially younger siblings experience when siblings move away, grow older and change. It was also chosen because it is relatively simple and can be fit into nearly any plot structure or length which would give greater flexibility later on as it was not clear when this scenario was created what form the entire game would take. Plus, it allowed explanation for more of the game elements such as levels being floors in the bank, and their being an explanation for why the player is having to modify the player character.

\section{Game Structure}\label{sec:cut/plannedLevels}

It was not expected that all of the following would get implemented in full outside of a best case scenario.

Whilst there are many plot structures we would have borrowed from, our earliest planning was to follow the three act structure. \autoref{tab:cutContent/tableOfLevels} shows the rough plan for this structure in terms of events, plot elements and programming features. Levels use the \texttt{<chapter/floor>--<room>} naming scheme (where chapters and floors can be used interchangably), mirroring the \texttt{<world>--<level>} scheme of mario bros.

We chose to think mainly in terms of the three act structure because we wanted a more episodic plot to allow for greater flexibility than more prescriptive structures such as \textquote{The Hero's journey}, or Save the Cat's \textquote{Beatsheet} \citep{SaveTheCat}. It could also be argued that this structure falls into a \textquote{Fichtean Curve} as described by \citet{plotStruct}. The main notable deviation from most definitions of the three act structure, including Duncan's, is that that Code Heist was going to start \textquote{In Medias Res}, with the Eve and the Player on the street opposite the bank.

It is common wisdom in game development that you introduce new mechanics in a safe environment before progressively increasing complexity and danger; before eventually combining it with other mechanics. This mirrors writing where you explore what you have already introduced before introducing more elements. This also mirrors how students are taught, starting with easier problems, before moving onto more difficult challengers. At time of writing, we have not been able to research the theory or structure behind this, but it is planned that the levels of the game will informally follow a similar progression to lessons in a school.

\begin{fullwidthtable}
    \centering
    \begin{tabular}{ccp{6.8cm}p{7.4cm}}\toprule
        Act & Floor & Thematic/Plot Content & Gameplay Content \\\midrule
        1
            & 0
                & Getting CatBot up to the bank and the ground floor of the bank, before the auto-controllers start getting jammed and breaking.
                & Short introduction to platforming using an auto-controlled CatBot.
                \\ \\
            & 1
                & CatBot starts to loose signal, plus it starts to become apparent that the bank might not be an ordinary bank.
                & Basic CatCode programming, recreating jumping, striking etc. before going to more complex scenarios. Starting out with just implementing sequences of functions executed on button presses before moving onto scripts ran every frame requiring the use of selection.
                \\ \\
        2
            & 2\&3
                & More complex security measures start to appear necessitating more complex programming scenarios. Narrative function shifts to tension between the Player and Eve as Eve is being secretive about the nature of what they are doing. Tensions reach a boiling point towards the end of floor 3.
                & More programming elements are introduced, alongside This is where variables, states and other puzzles described in \autoref{sec:cut/puzzleIdeas} will start to appear. Specifics cannot be provided until the difficulty of puzzles can be determined.
                \\ \\
        3
            & 4
                & Eve and the Player reconcile as she reveals that she is reveals that her motivation. This floor contains the hardest puzzles in the game before the final \textquote{boss}.
                & Continuation of the gameplay arc in floors 2 and 3, including all of the CatCode spec. However, hints will be less available as Eve is focusing on hacking into the Boss's account for the vault code.
                \\ \\
            & 5
                & The team reach the final floor on their mission, before facing off against the evil boss. It is revealed that the vault code is 1234, alongside an appropraite joke about insecure passwords. Once the boss is initially appears to be defeated, the self destruct is initiated alongside an escape sequence. 
                & This section would include all the gimmicks from previous floors, combining them. The boss would not be a traditional boss, instead being a scenario where the player is forced to program CatBot in safe areas before moving into more dangerous areas, whilst thinking ahead in terms of the vent mini-game.
                \\ \\
            & 6
                & This is the only part of the game not represented by a floor, instead it is an escape sequence throughout the bank as it self destructs. At the end the there will be a final phase to the final boss in the sky above the city, where the player programs the orbital death laser.
                & No new gameplay will be introduced here, rather covering everything before but with more (perceived) time constraints. This is meant to act more as a victory lap than a real challenge. Hacking the Orbital Death Laser is not intended to be difficult, but will be the first time a mechanic is not introduced in a safe environment.
                
    \\ \bottomrule\end{tabular}
    \caption[Outline of CodeHeist's story and gameplay]{Basic preliminary draft of CodeHeist in terms of plot and gameplay. Note: Story elements would be the first to be cut and were to be treated as asperational.}
    \label{tab:cutContent/tableOfLevels}
\end{fullwidthtable}

\section{Puzzle Ideas}\label{sec:cut/puzzleIdeas}

Below is a list of mechanics that could be implemented into the game for puzzles. It is not yet clear how usable or fun any of the following are.

\paragraph{Recreating jumping and movement} \textit{uses: selection, and functions}

\paragraph{Navigating vents} \textit{uses: definite iteration, functions, etc}. One of the earliest problems identified with the platformer medium was including iteration, a workaround was to create a mini game where the player programs a drone (potentially the cat's tail to make it an homage to the classic game snake) which would navigate the a series of vents using pre-programmed instructions as there is no signal to control it from the outside. At minimum this would allow the inclusion of iteration.

\paragraph{Hacking other robots} \textit{uses: various functions}. This would be a more complex system to implement, however, if we wanted the player to learn to read and modify existing code, this could be done through hacking various robots to break them, or make them less of a threat. This would be making a laser move slower or removing a robot's ability to attack for example.

\paragraph{Different States} \textit{uses: selection, and (potentially) variables}. To add further complexity to the CatBot controllers, we could add multiple different states that CatBot could enter which would vary how it behaves. For example, CatBot could enter a water state when in the water which would prevent jumping and require a swim function. This example would require the implementation of water. As a bonus, this teaches the player about states which are a concept common in programming, and may allow for an introduction to finite state machines.

\paragraph{Advanced movement} \textit{uses: various functions}. Similar to Celeste's movement \textquote{tech} such as \textquote{wave-dashing} \citep{celesteTech}, CodeHeist has various advanced techniques within it's movement. For example, you can strike immediately before jumping for extra distance. The physics engine used as a template for CatBot had even more of these. Because these can be hard to pull off with human hands, the player may be able to reimplement these using CatCode to occure in a single button press rather than multiple.

\paragraph{Limited editing} An extra layer of challenge that can be added on top of other scenarios is limiting where CatBot's code can be edited. This would increase the amount of forward thinking and planning the player would need to partake in to clear levels. It would also mean that players would have to think about how they bind keys as they may be limited in the number of keys they have access to, so may have to resort to combinations of keys.

\paragraph{Error-proof code} To add an extra layer of challenge to the game, we might be able to add optional objectives for creating error-proof code. This is code where, no mater what inputs you do in what state, CatBot can never throw errors. For example, the player has created code that wont allow the jump function to be called in midair.

Making this optional may also make the game more approachable as less experienced programmers will be able to make it through levels (using jump as an example) by avoiding pressing their jump key whilst in midair without implementing error prevention. This may, however, risk allowing players to avoid certain mechanics.

